\chapter{Esercizi sui diagrammi di stato}

Questo capitolo raccoglie gli esercizi della lezione dedicata ai diagrammi di stato. Per ognuno riportiamo il testo, il
ragionamento che porta alla soluzione, la soluzione proposta a lezione ridisegnata e un commento. In alcuni casi il commento
segnala un punto in cui la soluzione proposta \textbf{va corretta o completata}: sono proprio questi i dettagli su cui si
perdono punti all'esame, quindi vale la pena leggerli con attenzione.

\section{Un metodo per affrontare gli esercizi}

\Needspace{12\baselineskip}
Prima di disegnare conviene rispondere, nell'ordine, a quattro domande.
\begin{enumerate}
  \item \textbf{Quali sono gli stati?} Sono le situazioni in cui il sistema \textbf{reagisce in modo diverso} agli stessi eventi.
        Un buon indizio sono gli aggettivi e le espressioni del testo come «acceso», «in attesa», «in fase di stampa».
  \item \textbf{Quali sono gli eventi?} Sono le azioni dell'utente (premere un pulsante), i messaggi di altri sistemi (una
        richiesta via rete) e i \textbf{timer} («una volta al giorno», «ogni secondo»).
  \item \textbf{Ci sono stati che si raggruppano?} Se un evento ha lo stesso effetto da molti stati (ad esempio «spegni» da
        qualunque stato di funzionamento), quegli stati vanno racchiusi in uno \textbf{stato composito}, e la transizione parte
        dal suo bordo. Se due comportamenti procedono \textbf{indipendentemente}, servono \textbf{regioni parallele}.
  \item \textbf{Serve ricordare dove si era?} Se il testo dice «riprende», «torna come prima», «ripristina», serve uno
        pseudostato \textbf{history}.
\end{enumerate}

\info{{Contatori, guardie e azioni}}{
  Quando il testo parla di quantità (punti vita, secondi trascorsi, punteggi), quelle quantità \textbf{non} diventano stati: si
  tengono in variabili, che si modificano nelle \textbf{azioni} (\texttt{/ hp -= 5}) e si controllano nelle \textbf{guardie}
  (\texttt{[hp < 30]}). È la stessa idea del contatore del Capitolo 9.
}

% ============================================================
\section{Esercizio 1: una finestra del desktop}

\begin{esercizio}{La finestra}
Si descriva con uno statechart il comportamento di una generica finestra (ad esempio minimizzata, massimizzata, in modalità
finestra, \ldots) in un ambiente desktop basato su finestre, come quello di Microsoft Windows.

\textbf{Ragionamento.} Una finestra aperta può essere \textbf{visibile} oppure \textbf{ridotta a icona}. Quando è visibile, può
essere in modalità finestra o massimizzata. Quando la si riapre dalla barra delle applicazioni, deve tornare \textbf{come era
prima}: massimizzata se lo era. Serve quindi uno stato composito Visible con dentro i due sottostati, e uno pseudostato history.

\begin{center}
\begin{tikzpicture}
  \node[stato, minimum width=1.9cm, minimum height=3.6cm] (min) at (0,0) {};
  \node[font=\scriptsize\bfseries, anchor=north] at (min.north) {Minimized};
  \draw[rounded corners=8pt, line width=0.8pt, draw=q3!80!black, fill=q3!18] (2.9,-1.9) rectangle (8.0,1.9);
  \node[font=\scriptsize\bfseries] at (5.45,1.62) {Visible};
  \node[iniziale] (i0) at (5.45,2.75) {};
  \draw[tr] (i0) -- (5.45,1.9);
  \node[iniziale] (i1) at (4.0,0.75) {};
  \node[stato, fill=q3!35, minimum width=1.9cm] (win) at (6.2,0.75) {Windowed};
  \node[stato, fill=q3!35, minimum width=1.9cm] (max) at (6.2,-1.05) {Maximized};
  \draw[tr] (i1) -- (win);
  \draw[tr] ($(win.south)+(-0.4,0)$) -- node[left, font=\tiny]{maximize()} ($(max.north)+(-0.4,0)$);
  \draw[tr] ($(max.north)+(0.4,0)$) -- node[right, font=\tiny]{maximize()} ($(win.south)+(0.4,0)$);
  \node[circle, draw, fill=white, font=\tiny, inner sep=1pt, minimum size=0.36cm] (h) at (3.6,-1.05) {H};
  \draw[tr] (2.9,0.2) -- node[above, font=\tiny]{minimize()} (min.east |- 0,0.2);
  \draw[tr] (min.east |- 0,-1.05) -- node[above, font=\tiny]{taskbarClick()} (h);
  \node[finale] (f) at (9.4,0) {};
  \draw[tr] (8.0,0) -- node[above, font=\tiny]{close()} (f);
\end{tikzpicture}
\end{center}

\begin{itemize}
  \item All'apertura si entra in \textbf{Visible} e, dallo stato iniziale interno, in \textbf{Windowed}.
  \item Lo stesso evento \texttt{maximize()} fa passare da Windowed a Maximized e viceversa: è il pulsante che alterna le due
        modalità.
  \item \texttt{minimize()} parte dal \textbf{bordo} di Visible, quindi vale da entrambi i sottostati.
  \item \texttt{taskbarClick()} entra attraverso la \textbf{H}: la finestra torna massimizzata se lo era quando è stata ridotta a
        icona, in modalità finestra altrimenti.
\end{itemize}
\end{esercizio}

\warning{Che cosa si potrebbe migliorare}{
  Nel modello \texttt{close()} esce solo da Visible: una finestra ridotta a icona non si può chiudere. In Windows invece si può
  chiudere anche dalla barra delle applicazioni (tasto destro, \textit{Chiudi finestra}). Per modellarlo basta racchiudere
  Minimized e Visible in un altro stato composito (ad esempio \textbf{Open}) e far partire \texttt{close()} dal suo bordo.
}

% ============================================================
\section{Esercizio 2: l'orologio con la radio}

\begin{esercizio}{La radiosveglia}
Un orologio da tavolo, una volta acceso, mostra l'orario corrente sul proprio display LCD e, se l'utente preme un apposito
pulsante, può anche sintonizzarsi su stazioni radio e riprodurne le trasmissioni dalle casse integrate. Tramite un pulsante
«next station» è possibile passare alla stazione radio successiva, che verrà riprodotta dopo una breve fase di ricerca e
sintonizzazione.

\textbf{Ragionamento.} Quando l'orologio è acceso fa \textbf{due cose indipendenti}: mostra l'ora, sempre, e gestisce la radio,
che può essere accesa o spenta. Due comportamenti indipendenti si modellano con due \textbf{regioni parallele} dello stato On.
Dentro la radio accesa c'è a sua volta un piccolo ciclo: sintonizzazione, poi riproduzione.

\begin{center}
\begin{tikzpicture}
  \node[iniziale] (i) at (0,1.6) {};
  \node[stato, minimum width=1.6cm] (off) at (0,0) {Off};
  \draw[tr] (i) -- (off);
  \draw[rounded corners=8pt, line width=0.8pt, draw=q3!80!black, fill=q3!18] (2.1,-3.5) rectangle (11.4,2.2);
  \node[font=\scriptsize\bfseries] at (6.75,1.95) {On};
  \draw[dashed, line width=0.7pt] (2.1,0.75) -- (11.4,0.75);
  % regione 1
  \node[iniziale] (r1) at (2.8,1.3) {};
  \node[stato, fill=q3!35, minimum width=1.9cm] (st) at (5.4,1.3) {Show Time};
  \draw[tr] (r1) -- (st);
  % regione 2
  \node[iniziale] (r2) at (2.8,0.0) {};
  \node[stato, fill=q3!35, minimum width=1.7cm] (roff) at (4.3,0.0) {RadioOff};
  \draw[tr] (r2) -- (roff);
  \draw[rounded corners=8pt, line width=0.8pt, draw=q3!80!black, fill=q3!32] (6.4,-3.2) rectangle (11.0,0.45);
  \node[font=\scriptsize\bfseries] at (8.7,0.2) {RadioOn};
  \node[iniziale] (r3) at (7.0,-0.55) {};
  \node[stato, fill=q3!50, minimum width=2.2cm] (tun) at (9.0,-0.55) {Tuning};
  \node[stato, fill=q3!50, minimum width=2.2cm] (pl) at (9.0,-2.3) {Playing};
  \draw[tr] (r3) -- (tun);
  \draw[tr] ($(tun.south)+(0.45,0)$) -- node[right, font=\tiny]{done()} ($(pl.north)+(0.45,0)$);
  \draw[tr] ($(pl.north)+(-0.45,0)$) -- node[left, font=\tiny]{nextStation()} ($(tun.south)+(-0.45,0)$);
  \draw[tr] (roff.east) -- node[above, font=\tiny]{radio()} (6.4,0.0);
  \draw[tr] (6.4,-2.6) -- node[above, font=\tiny]{radio()} (4.3,-2.6) -- (roff.south);
  % accensione e spegnimento
  \draw[tr] (off.east) -- node[above, font=\tiny]{on()} (2.1,0);
  \draw[tr] (2.1,-3.1) -- node[above, font=\tiny]{off()} (0,-3.1) -- (off.south);
\end{tikzpicture}
\end{center}

\begin{itemize}
  \item Accendendo l'orologio (\texttt{on()}) si entra in On, e quindi in \textbf{entrambe} le regioni: si mostra l'ora e la radio è
        spenta.
  \item \texttt{radio()} accende la radio: si entra in RadioOn e subito in \textbf{Tuning}. Finita la sintonizzazione
        (\texttt{done()}) si passa in Playing; \texttt{nextStation()} riporta in Tuning per cercare la stazione successiva.
  \item Premendo di nuovo \texttt{radio()} si spegne la radio da qualunque sottostato di RadioOn (la transizione parte dal bordo),
        mentre l'ora continua a essere mostrata.
  \item \texttt{off()} parte dal bordo di On: spegne tutto, e termina entrambe le regioni.
\end{itemize}
\end{esercizio}

\info{Perché le regioni parallele}{
  Senza regioni parallele servirebbero stati come «ora + radio spenta», «ora + sintonizzazione», «ora + riproduzione»: il numero di
  stati si moltiplica. Con le regioni, i comportamenti indipendenti si descrivono separatamente e si \textbf{combinano} da soli. Qui
  la regione dell'ora è banale, ma rende esplicito che l'ora resta visibile qualunque cosa faccia la radio.
}

% ============================================================
\section{Esercizio 3: la schermata di login}

\begin{esercizio}{Il login}
La schermata di login di un'applicazione permette agli utenti di inserire le proprie credenziali ed accedere. Se il nome utente
inserito non è tra quelli presenti nel sistema, viene mostrato un warning dedicato. Altrimenti, se il nome è presente ma la password
è errata, viene mostrato un diverso warning e viene abilitato un pulsante per accedere alla funzionalità di reset password. Se le
credenziali sono corrette, si accede al sistema.

\textbf{Ragionamento.} Ogni «aspetto» diverso della schermata è uno stato: vuota, compilata, con il warning sul nome utente, con il
warning sulla password (e il pulsante di reset), accesso riuscito. Lo stesso evento \texttt{submit} porta in stati diversi a
seconda dei dati: si usano \textbf{guardie} mutuamente esclusive.

\begin{center}
\begin{tikzpicture}
  \node[iniziale] (i) at (0,1.2) {};
  \node[stato, minimum width=1.9cm] (empty) at (0,0) {LoginEmpty};
  \node[stato, minimum width=1.9cm] (fill) at (3.6,0) {LoginFilled};
  \node[stato, minimum width=2.7cm] (ue) at (10.0,1.4) {LoginUsernameError};
  \node[stato, minimum width=2.7cm] (pe) at (10.0,-1.4) {LoginPasswordError};
  \node[stato, minimum width=1.9cm] (ok) at (3.6,-3.1) {LoginSuccess};
  \node[finale] (f) at (7.6,-3.1) {};
  \draw[tr] (i) -- (empty);
  \draw[tr] (empty) -- node[above, font=\tiny]{fill(usr,pwd)} (fill);
  \draw[tr] ($(fill.north)+(0.5,0)$) |- node[etich, pos=0.75]{submit(usr,pwd) [!userExists(usr)]} (ue.west);
  \draw[tr] (ue.north) |- ++(0,0.45) -| node[etich, pos=0.25]{fill(usr,pwd)} ($(fill.north)+(-0.5,0)$);
  \draw[tr] ($(fill.east)+(0,0.12)$) -| node[etich, pos=0.3, text width=4.2cm]{submit(usr,pwd) [userExists(usr) \&\\ !checkLogin(usr,pwd)]} ($(pe.north)+(-0.6,0)$);
  \draw[tr] (pe.west) -- node[etich, pos=0.5]{fill(usr,pwd)} ($(fill.south east)+(0,-0.05)$ |- pe.west) -- ($(fill.south east)+(-0.25,0)$);
  \draw[tr] (fill.south) -- node[etich, text width=4.2cm]{submit(usr,pwd) [userExists(usr) \&\\ checkLogin(usr,pwd)]} (ok.north);
  \draw[tr] (ok) -- node[above, font=\tiny]{redirectTimeout()} (f);
  % aggiunta
  \node[stato, minimum width=2.4cm, draw=swred, fill=swred!8, dashed] (rp) at (10.0,-3.6) {ResetPassword};
  \draw[tr, draw=swred, dashed] (pe.south) -- node[right, font=\tiny, text=swred]{resetPassword()} (rp.north);
\end{tikzpicture}
\end{center}

\begin{itemize}
  \item Compilando i campi si passa da LoginEmpty a LoginFilled. Da qui \texttt{submit} porta in \textbf{tre} stati diversi, e le
        tre guardie coprono tutti i casi senza sovrapporsi: utente inesistente; utente esistente con password errata; utente
        esistente con password corretta.
  \item Dai due stati di errore, ricompilando i campi (\texttt{fill}) si torna in LoginFilled e si può riprovare.
  \item Dopo l'accesso riuscito, un timer (\texttt{redirectTimeout()}) porta fuori dalla schermata di login.
\end{itemize}
\end{esercizio}

\warning{Che cosa va aggiunto alla soluzione proposta}{
  Il testo dice che con la password errata «viene abilitato un pulsante per accedere alla funzionalità di reset password», ma la
  soluzione proposta a lezione non modella mai la pressione di quel pulsante: il pulsante esiste, ma premerlo non fa nulla. Nel
  disegno qui sopra è aggiunta in rosso la transizione \texttt{resetPassword()} verso uno stato \textbf{ResetPassword}, che esce
  solo da LoginPasswordError: è proprio lì che il pulsante è abilitato.

  La nota della slide segnala un altro limite, voluto: in questo modello l'utente non può chiudere la schermata prima di aver
  completato il login, perché l'unico stato finale si raggiunge da LoginSuccess.
}

% ============================================================
\section{Esercizio 4: la stampante}

\begin{esercizio}{La stampante di rete}
Una stampante, previa accensione, resta in attesa di ricevere via rete documenti da stampare. In presenza di richieste, la
stampante procede alla stampa. Quando è accesa e non è in fase di stampa, la stampante, una volta al giorno, effettua la pulizia
delle testine. Inoltre, sempre con cadenza giornaliera, la stampante scarica e installa aggiornamenti dalla casa madre. In questo
caso, la stampante interrompe qualsiasi attività in corso per effettuare l'aggiornamento, e le riprende ad aggiornamento
effettuato.

\textbf{Ragionamento.} Le frasi chiave sono tre.
\begin{itemize}
  \item «Quando non è in fase di stampa \ldots\ effettua la pulizia»: la pulizia parte solo dallo stato di attesa (Idle), non da
        Printing.
  \item «Interrompe \textbf{qualsiasi} attività»: l'aggiornamento parte dal bordo di uno stato composito che racchiude tutte le
        attività normali (Available).
  \item «Le \textbf{riprende}»: al ritorno serve una \textbf{history}.
\end{itemize}

\begin{center}
\begin{tikzpicture}
  \node[iniziale] (i) at (0,1.7) {};
  \node[stato, minimum width=1.4cm] (off) at (0,0.6) {Off};
  \draw[tr] (i) -- (off);
  \draw[rounded corners=8pt, line width=0.8pt, draw=q3!80!black, fill=q3!15] (1.8,-4.6) rectangle (13.3,2.0);
  \node[font=\scriptsize\bfseries] at (7.55,1.75) {On};
  \node[iniziale] (ion) at (2.3,0.6) {};
  % Available
  \draw[rounded corners=8pt, line width=0.8pt, draw=q3!80!black, fill=q3!28] (2.9,-4.3) rectangle (8.2,1.4);
  \node[font=\scriptsize\bfseries] at (5.55,1.15) {Available};
  \draw[tr] (ion) -- (2.9,0.6);
  \node[iniziale] (ia) at (3.3,0.6) {};
  \node[stato, fill=q3!45, minimum width=1.5cm] (idle) at (4.0,-0.5) {Idle};
  \node[stato, fill=q3!45, minimum width=1.5cm] (pr) at (7.1,-0.5) {Printing};
  \node[stato, fill=q3!45, minimum width=1.5cm] (cl) at (4.0,-2.9) {Cleaning};
  \draw[tr] (ia) -- (idle);
  \draw[tr] ($(idle.east)+(0,0.15)$) -- node[etich, above]{request()} ($(pr.west)+(0,0.15)$);
  \draw[tr] ($(pr.west)+(0,-0.15)$) -- node[etich, below]{done()} ($(idle.east)+(0,-0.15)$);
  \draw[tr] ($(idle.south)+(0.3,0)$) -- node[etich, right]{cleaningTimeout()} ($(cl.north)+(0.3,0)$);
  \draw[tr] ($(cl.north)+(-0.3,0)$) -- node[etich, left]{cleaningDone()} ($(idle.south)+(-0.3,0)$);
  \node[circle, draw, fill=white, font=\tiny, inner sep=1pt, minimum size=0.36cm] (h) at (7.3,-3.6) {H};
  % Updating
  \draw[rounded corners=8pt, line width=0.8pt, draw=q3!80!black, fill=q3!28] (10.0,-4.3) rectangle (12.9,1.4);
  \node[font=\scriptsize\bfseries] at (11.45,1.15) {Updating};
  \node[iniziale] (iu) at (11.45,0.7) {};
  \node[stato, fill=q3!45, minimum width=2.0cm] (dl) at (11.45,-0.3) {Downloading};
  \node[stato, fill=q3!45, minimum width=2.0cm] (ins) at (11.45,-1.9) {Installing};
  \node[finale] (fu) at (11.45,-3.2) {};
  \draw[tr] (iu) -- (dl);
  \draw[tr] (dl) -- node[right, font=\tiny]{completed()} (ins);
  \draw[tr] (ins) -- node[right, font=\tiny]{installed()} (fu);
  \draw[tr] (8.2,0.6) -- node[above, font=\tiny]{updateTimeout()} (10.0,0.6);
  \draw[tr] (10.0,-3.6) -- node[above, font=\tiny]{updateDone()} (h);
  \draw[tr] (1.8,-3.9) -- node[above, font=\tiny]{off()} (0,-3.9) -- (off.south);
  \draw[tr] (off.east) -- node[above, font=\tiny]{on()} (1.8,0.6);
\end{tikzpicture}
\end{center}

\begin{itemize}
  \item Accesa la stampante si entra in On, quindi in \textbf{Available}, quindi in \textbf{Idle}.
  \item Da Idle una richiesta avvia la stampa; un timer giornaliero avvia la pulizia. Da Printing la pulizia non può partire: la
        transizione \texttt{cleaningTimeout()} esce solo da Idle.
  \item Il timer degli aggiornamenti esce dal \textbf{bordo} di Available: interrompe stampa, pulizia o attesa, qualunque cosa
        la stampante stesse facendo.
  \item Finito l'aggiornamento si torna in Available attraverso la \textbf{H}: la stampante riprende l'attività interrotta.
\end{itemize}
\end{esercizio}

\warning{Che cosa si potrebbe migliorare}{
  \begin{itemize}
    \item Updating termina con il suo stato finale interno (dopo \texttt{installed()}), ma per uscirne la soluzione usa un evento
          in più, \texttt{updateDone()}. È più coerente usare una \textbf{transizione di completamento}: una transizione
          \textbf{senza trigger} che parte dal bordo di Updating e scatta da sola quando lo stato composito raggiunge il suo stato
          finale. Così non serve inventare un evento che il testo non menziona.
    \item Il modello non dice che cosa succede alle richieste di stampa che arrivano \textbf{durante} l'aggiornamento: in Updating
          l'evento \texttt{request()} non ha transizioni, quindi viene ignorato e il documento va perso. Una stampante reale le
          metterebbe in coda: va almeno chiarito con il committente.
  \end{itemize}
}

\info{Shallow o deep?}{
  Qui basta la history \textbf{shallow} (H) perché Idle, Printing e Cleaning sono stati semplici. Se Printing avesse a sua volta dei
  sottostati (ad esempio «riscaldamento» e «stampa pagina») e volessimo riprendere esattamente dal sottostato interrotto, servirebbe
  la history \textbf{deep} (H*).
}

% ============================================================
\section{Esercizio 5: il malware}

\begin{esercizio}{Il malware}
Un malware, una volta installato su un PC, rimane latente fino a quando l'utente non apre Internet Explorer. A quel punto, il malware
si attiva e, in parallelo, ricerca informazioni sensibili nei dischi rigidi e nella memoria del PC. Terminate queste attività, il
malware sfrutta una vulnerabilità di Internet Explorer per inviare le informazioni raccolte a un server remoto. Se Internet Explorer
viene chiuso prima dell'invio delle informazioni, il malware salva le informazioni trovate e riprova ad inviarle al successivo avvio
di Internet Explorer.

\textbf{Ragionamento.} «In parallelo» indica due \textbf{regioni parallele}. «Terminate queste attività» indica che si esce dallo
stato composito solo quando \textbf{entrambe} le regioni hanno finito, cioè una \textbf{transizione di completamento}. A quel punto
bisogna sapere se Internet Explorer è ancora aperto: lo decide una \textbf{guardia}.

\begin{center}
\begin{tikzpicture}
  \node[iniziale] (i) at (0,1.2) {};
  \node[stato, minimum width=1.4cm] (idle) at (0,0) {Idle};
  \draw[tr] (i) -- (idle);
  \draw[rounded corners=8pt, line width=0.8pt, draw=q3!80!black, fill=q3!18] (2.4,-2.0) rectangle (7.6,2.0);
  \node[font=\scriptsize\bfseries] at (5.0,1.75) {Active};
  \draw[dashed, line width=0.7pt] (2.4,0) -- (7.6,0);
  \node[iniziale] (a1) at (2.9,0.9) {};
  \node[stato, fill=q3!35, minimum width=2.4cm] (sm) at (4.7,0.9) {SearchingMemory};
  \node[finale] (f1) at (7.1,0.9) {};
  \draw[tr] (a1) -- (sm); \draw[tr] (sm) -- node[above, font=\tiny]{done()} (f1);
  \node[iniziale] (a2) at (2.9,-1.0) {};
  \node[stato, fill=q3!35, minimum width=2.4cm] (sd) at (4.7,-1.0) {SearchingDisks};
  \node[finale] (f2) at (7.1,-1.0) {};
  \draw[tr] (a2) -- (sd); \draw[tr] (sd) -- node[above, font=\tiny]{done()} (f2);
  \draw[tr] (idle) -- node[above, font=\tiny]{IEOpened()} (2.4,0);
  % dopo Active
  \node[stato, minimum width=2.4cm, align=left] (up) at (11.9,1.3) {\makebox[2.4cm]{\textbf{UploadingData}}\\[-2pt]\rule{2.4cm}{0.4pt}\\do / upload()};
  \node[stato, minimum width=2.4cm] (sw) at (10.5,-2.1) {SaveDataAndWait};
  \node[finale] (fe) at (15.3,1.3) {};
  \draw[tr] (7.6,1.3) -- node[above, font=\tiny]{[dataAcquired \& IEOpen]} (up.west);
  \draw[tr] (5.0,-2.0) |- node[etich, pos=0.7, above]{[dataAcquired \& !IEOpen]} (sw.west);
  \draw[tr] ($(up.south)+(-0.7,0)$) -- node[etich, left, pos=0.4]{IEClosed()} ($(sw.north)+(0.7,0)$);
  \draw[tr] (sw.east) -| node[etich, right, pos=0.75]{IEOpened()} ($(up.south)+(0.6,0)$);
  \draw[tr] (up.east) -- node[etich, above]{uploadDone()} (fe);
\end{tikzpicture}
\end{center}

\begin{itemize}
  \item All'apertura di Internet Explorer il malware si attiva ed entra in \textbf{Active}: le due ricerche partono insieme, una per
        regione. Ciascuna termina nel proprio stato finale.
  \item Quando \textbf{tutte e due} le regioni sono nel loro stato finale, Active è completato e scatta una delle due transizioni
        senza trigger, a seconda della guardia: se Internet Explorer è ancora aperto si inviano i dati, altrimenti li si salva e si
        aspetta.
  \item Se Internet Explorer viene chiuso durante l'invio si passa in SaveDataAndWait; alla riapertura si riprova l'invio.
\end{itemize}
\end{esercizio}

\info{Le transizioni di completamento}{
  Le due frecce che escono da Active non hanno trigger: hanno solo una guardia. Sono \textbf{transizioni di completamento}: scattano
  quando lo stato da cui partono ha \textbf{finito il suo lavoro}, cioè, per uno stato composito, quando ogni sua regione ha
  raggiunto il proprio stato finale. Le due guardie sono complementari, quindi ne scatta sempre esattamente una.
}

% ============================================================
\section{Esercizio 6: il giocatore di un videogioco}

\begin{esercizio}{Il giocatore}
In un videogame, il giocatore inizia la partita con una salute pari a 100 HP. Durante la partita, il giocatore può subire attacchi
diretti, che diminuiscono la salute residua di 5 HP. Inoltre, il giocatore può essere avvelenato. L'avvelenamento prevede una fase
iniziale che dura 5 secondi, in cui il giocatore subisce 2 HP di danno al secondo, e una fase acuta, che dura 3 secondi e durante la
quale il giocatore subisce 10 HP di danno al secondo. Mentre il giocatore è avvelenato in fase acuta, i danni inflitti da attacchi
diretti raddoppiano. Quando la salute scende al di sotto della soglia critica di 30 HP, il giocatore passa in modalità «berserk».
Quando è in questa modalità, il giocatore è immune all'avvelenamento e si cura di 5 HP al secondo, ma subisce danni doppi dai colpi
diretti. Quando i punti salute scendono a zero, il giocatore muore e la partita termina.

\textbf{Ragionamento.} Le quantità (salute, secondi trascorsi) vanno in due variabili, \texttt{hp} e \texttt{counter}. Gli stati sono
le situazioni in cui il giocatore reagisce in modo diverso: sano, avvelenato (in due fasi), berserk. Il passare del tempo si
modella con un evento \texttt{timeout1sec()} che arriva ogni secondo. La morte può avvenire da qualunque stato: tutto il gioco sta
dentro uno stato composito Alive, da cui si esce con una guardia sulla salute.

\begin{center}
\begin{tikzpicture}[lab/.style={font=\tiny, inner sep=1pt}]
  \node[iniziale] (i) at (0.5,1.0) {};
  \draw[rounded corners=8pt, line width=0.8pt, draw=q3!80!black, fill=q3!15] (0,-6.0) rectangle (15.4,0.5);
  \node[font=\scriptsize\bfseries] at (7.7,0.25) {Alive};
  \node[font=\tiny, anchor=west] at (0.15,0.0) {entry / hp = 100};
  \draw[tr] (i) -- (0.5,0.5);
  % Healthy e Berserk
  \node[iniziale] (ih) at (0.6,-1.4) {};
  \node[stato, fill=q3!35, minimum width=1.8cm] (he) at (2.6,-1.4) {Healthy};
  \draw[tr] (ih) -- (he);
  \path (he) edge[tr, loop above, out=110, in=70, min distance=8mm] node[lab, above]{directHit() / hp -= 5} (he);
  \node[stato, fill=q3!35, minimum width=1.8cm] (be) at (2.6,-4.2) {Berserk};
  \draw[tr] ($(he.south)+(0.4,0)$) -- node[lab, right]{[hp < 30]} ($(be.north)+(0.4,0)$);
  \draw[tr] ($(be.north)+(-0.4,0)$) -- node[lab, left]{[hp >= 30]} ($(he.south)+(-0.4,0)$);
  \path (be) edge[tr, loop left, min distance=8mm] node[lab, left, align=right]{timeout1sec()\\/ hp += 5} (be);
  \path (be) edge[tr, loop below, out=-110, in=-70, min distance=8mm] node[lab, below]{directHit() / hp -= 10} (be);
  % Poisoned
  \draw[rounded corners=8pt, line width=0.8pt, draw=q3!80!black, fill=q3!30] (5.6,-5.0) rectangle (15.1,-0.6);
  \node[font=\scriptsize\bfseries] at (10.35,-0.85) {Poisoned};
  \node[font=\tiny, anchor=west] at (5.75,-1.15) {entry / counter = 0};
  \draw[tr] (he.east) -- node[lab, above]{poisoned()} (5.6,-1.4);
  \node[iniziale] (ip) at (9.6,-1.5) {};
  \node[stato, fill=q3!50, minimum width=2.0cm] (ini) at (9.6,-2.4) {InitialPhase};
  \node[stato, fill=q3!50, minimum width=2.0cm] (acu) at (9.6,-4.1) {AcutePhase};
  \draw[tr] (ip) -- (ini);
  \draw[tr] (ini) -- node[lab, right]{[counter == 5]} (acu);
  \path (ini) edge[tr, loop right, min distance=8mm] node[lab, right, align=left]{timeout1sec() [counter < 5]\\/ hp -= 2; counter++} (ini);
  \path (acu) edge[tr, loop right, min distance=8mm] node[lab, right, align=left]{timeout1sec() [counter < 8]\\/ hp -= 10; counter++} (acu);
  \path (ini) edge[tr, loop left, min distance=8mm] node[lab, left]{directHit() / hp -= 5} (ini);
  \path (acu) edge[tr, loop left, min distance=8mm] node[lab, left]{directHit() / hp -= 10} (acu);
  \draw[tr] (5.6,-3.6) -- node[lab, above]{[hp < 30]} (be.east |- 0,-3.6) ;
  \draw[tr] ($(acu.south)+(0,0)$) -- (9.6,-5.5) -- node[lab, above]{[counter == 8]} (1.6,-5.5) -- (1.6,-1.4) -- (he.west);
  % morte
  \node[finale] (f) at (7.7,-6.9) {};
  \draw[tr] (7.7,-6.0) -- node[lab, right]{[hp <= 0]} (f);
\end{tikzpicture}
\end{center}

\begin{itemize}
  \item Entrando in Alive la salute vale 100 (\texttt{entry}). Da sano, ogni colpo toglie 5 HP.
  \item L'avvelenamento porta in \textbf{Poisoned}, che azzera il contatore dei secondi (\texttt{entry}). Nella fase iniziale ogni
        secondo toglie 2 HP; dopo 5 secondi si passa alla fase acuta, dove ogni secondo toglie 10 HP e i colpi diretti valgono
        doppio. Dopo altri 3 secondi (contatore a 8) l'avvelenamento finisce e si torna sani.
  \item Sotto i 30 HP si passa in \textbf{Berserk}, anche durante l'avvelenamento: la transizione parte dal bordo di Poisoned e lo
        interrompe. In Berserk non c'è la transizione \texttt{poisoned()}: è così che si modella l'immunità.
  \item Quando la salute arriva a zero si esce da Alive, da qualunque stato.
\end{itemize}
\end{esercizio}

\Needspace{24\baselineskip}
\error{Che cosa va corretto nella soluzione proposta}{
  Il disegno qui sopra contiene già tre correzioni rispetto alla soluzione mostrata a lezione. Vale la pena vederle una per una,
  perché sono errori tipici.
  \begin{itemize}
    \item \textbf{La fase acuta dura un secondo di troppo.} La soluzione proposta usa \texttt{[counter <= 8]} sull'autotransizione
          della fase acuta e \texttt{[counter > 8]} per uscire. Entrando in fase acuta \texttt{counter} vale 5: i tick avvengono
          con \texttt{counter} pari a 5, 6, 7 \textbf{e 8}, cioè quattro secondi e 40 HP di danno invece di 30. Le guardie corrette
          sono \texttt{[counter < 8]} e \texttt{[counter == 8]}.
    \item \textbf{Nella fase iniziale le guardie si sovrappongono.} Con \texttt{[counter <= 5]} sull'autotransizione e
          \texttt{[counter = 5]} sul passaggio alla fase acuta, quando \texttt{counter} vale 5 sono vere \textbf{entrambe} le
          guardie. Il modello funziona solo grazie a una regola implicita: la transizione senza evento scatta subito, prima che
          arrivi il tick successivo. Con \texttt{[counter < 5]} le guardie diventano mutuamente esclusive e il modello non dipende
          più da questo dettaglio.
    \item \textbf{Il giocatore potrebbe non morire mai.} La soluzione proposta esce da Alive con \texttt{[hp = 0]}. Ma la salute
          scende a passi di 2, 5 o 10: con 3 HP, un colpo diretto la porta a $-2$, e la guardia \texttt{hp = 0} non sarà mai vera.
          La guardia corretta è \texttt{[hp <= 0]}.
  \end{itemize}
}

\info{Una lezione generale sulle guardie con i contatori}{
  Quando un contatore deve contare $n$ eventi, conviene verificare a mano i valori estremi: con che valore entra, con quali valori
  scatta l'autotransizione, con quale valore si esce. E le guardie su una quantità che varia «a salti» vanno scritte con
  \texttt{<=} o \texttt{>=}, mai con \texttt{==}, a meno di essere certi che il valore esatto venga raggiunto.
}

% ============================================================
\section{Esercizio 7: un game di tennis}

\begin{esercizio}{Il punteggio di un game}
Si modelli l'andamento del punteggio di un game di una partita di tennis.

\textbf{Ragionamento.} Ricordiamo le regole. Ogni giocatore parte da 0 e, a ogni punto vinto, passa a 15, 30, 40. Con 40, vincendo
un altro punto vince il game, a meno che anche l'avversario sia a 40: in quel caso si è in \textbf{parità} (\textit{deuce}). Dalla
parità, chi vince il punto va in \textbf{vantaggio}; se vince anche il punto successivo vince il game, altrimenti si torna in parità.
I punteggi da 0 a 40 si tengono in due variabili \texttt{P1} e \texttt{P2}; parità e vantaggio, invece, cambiano il modo in cui si
reagisce ai punti, quindi sono \textbf{stati}. L'evento \texttt{P1*} indica «punto vinto dal giocatore 1».

\begin{center}
\begin{tikzpicture}[lab/.style={font=\tiny, inner sep=1pt, fill=white}]
  \node[iniziale] (i) at (0,1.0) {};
  \node[stato, minimum width=2.2cm] (pg) at (0,0) {Punto in gioco};
  \draw[tr] (i) -- node[lab, right]{/ P1 = 0; P2 = 0} (pg);
  \path (pg) edge[tr, loop left, min distance=10mm] node[lab, left, align=right]{P1* [P1 < 30 || (P1 == 30 \&\& P2 < 40)]\\/ P1 = prossimo(P1)} (pg);
  \path (pg) edge[tr, loop right, min distance=10mm] node[lab, right, align=left]{P2* [P2 < 30 || (P2 == 30 \&\& P1 < 40)]\\/ P2 = prossimo(P2)} (pg);
  \node[stato, minimum width=1.6cm] (de) at (0,-2.6) {Deuce};
  \draw[tr] ($(pg.south)+(-0.5,0)$) -- node[lab, left, align=right]{P1* [P1 == 30 \&\& P2 == 40]} ($(de.north)+(-0.5,0)$);
  \draw[tr] ($(pg.south)+(0.5,0)$) -- node[lab, right, align=left]{P2* [P2 == 30 \&\& P1 == 40]} ($(de.north)+(0.5,0)$);
  \node[stato, minimum width=1.9cm] (v1) at (-2.6,-4.4) {Vantaggio P1};
  \node[stato, minimum width=1.9cm] (v2) at (2.6,-4.4) {Vantaggio P2};
  \draw[tr] ($(de.south west)+(0.2,0)$) -- node[lab, left]{P1*} ($(v1.north)+(0.3,0)$);
  \draw[tr] ($(v1.north)+(0.8,0)$) -- node[lab, right]{P2*} ($(de.south west)+(0.55,0)$);
  \draw[tr] ($(de.south east)+(-0.2,0)$) -- node[lab, right]{P2*} ($(v2.north)+(-0.3,0)$);
  \draw[tr] ($(v2.north)+(-0.8,0)$) -- node[lab, left]{P1*} ($(de.south east)+(-0.55,0)$);
  \node[stato, minimum width=1.9cm] (fg) at (0,-6.3) {Fine game};
  \node[finale] (f) at (0,-7.4) {};
  \draw[tr] (fg) -- (f);
  \draw[tr] (v1.south) |- node[lab, pos=0.7]{P1* / vince P1} ($(fg.west)+(0,0.15)$);
  \draw[tr] (v2.south) |- node[lab, pos=0.7]{P2* / vince P2} ($(fg.east)+(0,0.15)$);
  \draw[tr] (pg.west) ++(0,-0.25) -- ++(-3.8,0) |- node[lab, pos=0.25, align=center]{P1* [P1 == 40 \&\& P2 < 40]\\/ vince P1} ($(fg.west)+(0,-0.15)$);
  \draw[tr] (pg.east) ++(0,-0.25) -- ++(3.8,0) |- node[lab, pos=0.25, align=center]{P2* [P2 == 40 \&\& P1 < 40]\\/ vince P2} ($(fg.east)+(0,-0.15)$);
\end{tikzpicture}
\end{center}

\begin{itemize}
  \item In \textbf{Punto in gioco} ogni punto fa avanzare il punteggio di chi lo vince (0, 15, 30, 40), finché non si verifica uno
        dei casi speciali.
  \item Se un giocatore a 40 vince un punto mentre l'altro è sotto i 40, vince il game.
  \item Se un giocatore a 30 vince il punto mentre l'altro è già a 40, si arriva a 40 pari: \textbf{Deuce}.
  \item Da Deuce si va in vantaggio; dal vantaggio si vince il game oppure si torna in parità.
\end{itemize}
\end{esercizio}

\error{Che cosa va corretto nella soluzione proposta}{
  \begin{itemize}
    \item \textbf{Guardie sovrapposte che portano a un vicolo cieco.} La soluzione proposta usa per l'autotransizione la guardia
          \texttt{[P1 : 0, 15, 30]} («P1 vale 0, 15 o 30»). Ma quando P1 vale 30 e P2 vale 40, è vera \textbf{anche} la guardia
          della transizione verso Deuce. Se scattasse l'autotransizione, P1 diventerebbe 40 restando in Punto in gioco, con il
          punteggio 40--40: da lì nessuna transizione è più attivabile correttamente (la vittoria richiede l'avversario sotto i 40,
          Deuce richiede P1 a 30), e il game non finisce più. Nel disegno l'autotransizione esclude il caso
          \texttt{P1 == 30 \&\& P2 == 40}, e lo stesso vale, simmetricamente, per P2.
    \item \textbf{Manca lo stato finale}, e lo stato iniziale è sostituito da uno stato «Inizio partita» con un evento «Inizio game».
          Per un game basta uno pseudostato iniziale (che azzera i punteggi) e uno stato finale dopo Fine game.
    \item L'azione \texttt{P1++} va intesa come «passa al punteggio successivo» (da 30 si va a 40, non a 31): nel disegno l'abbiamo
          scritta \texttt{prossimo(P1)} per non lasciare dubbi.
  \end{itemize}
}

% ============================================================
\section{Esercizio 8: una partita di tennis}

\begin{esercizio}{La partita intera}
Si modelli l'andamento del punteggio di un'intera partita di tennis (tre set su cinque).

\textbf{Ragionamento.} Un set si vince arrivando ad almeno 6 game con almeno 2 game di vantaggio (6--4, 7--5), oppure 7--6. La
partita si vince arrivando a 3 set. La soluzione proposta a lezione estende il diagramma del game aggiungendo, \textbf{sullo stesso
livello}, gli stati \textit{Game vinto P1/P2} e \textit{Set vinto P1/P2}, con contatori \texttt{GameP1}, \texttt{GameP2},
\texttt{SetP1}, \texttt{SetP2}. Il risultato funziona, ma ha oltre venti transizioni con guardie lunghe che si incrociano, ed è
molto difficile da leggere e da verificare.

Gli stati compositi permettono una soluzione molto più leggibile, organizzata su tre livelli come la partita stessa: la
partita contiene i set, il set contiene i game, e il game è \textbf{esattamente} il diagramma dell'esercizio precedente.

\begin{center}
\begin{tikzpicture}[lab/.style={font=\tiny, inner sep=1pt, fill=white, align=center}]
  \node[iniziale] (i) at (0.4,0.9) {};
  \draw[rounded corners=8pt, line width=0.8pt, draw=q3!80!black, fill=q3!12] (0,-5.6) rectangle (13.0,0.4);
  \node[font=\scriptsize\bfseries] at (6.5,0.15) {Partita};
  \node[font=\tiny, anchor=west] at (0.15,-0.15) {entry / SetP1 = 0; SetP2 = 0};
  \draw[tr] (i) -- (0.4,0.4);
  % Set
  \draw[rounded corners=8pt, line width=0.8pt, draw=q3!80!black, fill=q3!25] (0.4,-4.6) rectangle (9.2,-0.5);
  \node[font=\scriptsize\bfseries] at (4.8,-0.75) {Set};
  \node[font=\tiny, anchor=west] at (0.55,-1.05) {entry / GameP1 = 0; GameP2 = 0};
  \node[iniziale] (is) at (1.2,-2.3) {};
  \node[stato, fill=q3!45, minimum width=2.6cm, minimum height=1.1cm, align=center] (g) at (3.4,-2.3)
    {\textbf{Game}\\[-1pt]{\tiny (il diagramma}\\[-3pt]{\tiny dell'Esercizio 7)}};
  \draw[tr] (is) -- (g);
  \node[decisione] (dg) at (6.6,-2.3) {};
  \draw[tr] (g) -- node[lab, above]{/ GameX++} (dg);
  \node[font=\tiny, text=black!60] at (6.6,-1.85) {vincitore: X};
  \draw[tr] (dg.south) |- node[lab, pos=0.75]{[set non finito]} ($(g.south)+(0,-0.9)$) -- (g.south);
  \node[finale] (fs) at (8.6,-2.3) {};
  \draw[tr] (dg) -- node[lab, above]{[set finito]} (fs);
  % dopo il set
  \node[decisione] (dp) at (10.6,-2.3) {};
  \draw[tr] (9.2,-2.3) -- node[lab, above]{/ SetX++} (dp);
  \draw[tr] (dp.south) |- node[lab, pos=0.75]{[SetX < 3]} (4.8,-5.1) -- (4.8,-4.6);
  \node[finale] (fp) at (12.4,-2.3) {};
  \draw[tr] (dp) -- node[lab, above]{[SetX == 3]} (fp);
  \node[finale] (fe) at (6.5,-6.5) {};
  \draw[tr] (6.5,-5.6) -- (fe);
\end{tikzpicture}
\end{center}

\begin{itemize}
  \item \textbf{Game} è uno stato composito che contiene, tale e quale, il diagramma del game: quando il game finisce (stato
        finale interno), una transizione di completamento incrementa i game vinti da chi ha vinto e arriva a un
        \textbf{rombo}.
  \item Il rombo è lo pseudostato di \textbf{scelta} (\textit{choice}): come nei diagrammi di attività, sceglie un'uscita in base
        alle guardie. Se il set non è finito si gioca un altro game (rientrando in Game, i punteggi ripartono da zero); altrimenti
        il set termina.
  \item Allo stesso modo, finito il set, si incrementano i set vinti e si decide se giocare un altro set (rientrando in Set, i
        game ripartono da zero) o chiudere la partita.
  \item La condizione «set finito» si scrive una volta sola: \texttt{(GameX >= 6 \&\& GameX - GameY >= 2) || GameX == 7}, dove X è
        chi ha appena vinto il game e Y l'avversario.
\end{itemize}
\end{esercizio}

\info{{Cosa manca, in entrambe le soluzioni}}{
  Con le regole moderne, sul 6--6 non si gioca un game normale ma un \textbf{tie-break}, che ha un punteggio diverso (1, 2, 3, \ldots,
  fino ad almeno 7 con 2 punti di vantaggio). La soluzione proposta lo ignora e considera semplicemente il 7--6. Nel modello
  gerarchico aggiungerlo è facile: basta un secondo stato composito \textbf{TieBreak} accanto a Game, in cui si entra dal rombo con la
  guardia \texttt{[GameP1 == 6 \&\& GameP2 == 6]}. Nel modello piatto, invece, servirebbero molte altre transizioni.
}

\warning{Quando usare gli stati compositi}{
  Il confronto fra le due soluzioni dell'Esercizio 8 è il motivo per cui esistono gli stati compositi: un diagramma piatto con
  decine di transizioni incrociate è formalmente corretto, ma nessuno riesce a verificarlo. Se nel testo compare una struttura
  «annidata» (una partita fatta di set, fatti di game, fatti di punti), il diagramma dovrebbe rispecchiarla.
}

% ============================================================
\gbox{In sintesi}{azure-gradient-6!80!black}{
\begin{itemize}
  \item Gli \textbf{stati} sono le situazioni in cui il sistema reagisce in modo diverso agli stessi eventi; le \textbf{quantità}
        (salute, tempo, punteggi) vanno in variabili, aggiornate nelle azioni e controllate nelle guardie.
  \item Un evento che vale «da qualunque stato» suggerisce uno \textbf{stato composito} con la transizione sul bordo; comportamenti
        indipendenti suggeriscono \textbf{regioni parallele}; «riprende» o «torna come prima» suggeriscono una \textbf{history}.
  \item Una transizione \textbf{senza trigger} che esce da uno stato composito scatta quando lo stato \textbf{completa} il suo lavoro
        (tutte le regioni nel loro stato finale).
  \item Le guardie delle transizioni con lo stesso evento devono essere \textbf{mutuamente esclusive}; con i contatori vanno
        verificati a mano i valori estremi, e le soglie raggiunte «a salti» vanno scritte con \texttt{<=} o \texttt{>=}.
  \item Per strutture annidate (partita, set, game) gli stati compositi e lo pseudostato di \textbf{scelta} (rombo) producono
        diagrammi molto più leggibili di un grafo piatto.
\end{itemize}
}
